% =====================================================================
%  ch20-unconditional-i1.tex  --  the strong-coupling window at [T]
%  Sources (all read for this chapter): research/unconditional/
%  {PROOF-P3-UNCONDITIONAL, P3-FOURWALL-DERIVATION, P3-DRESSED-MASS,
%   P3-CYLINDER-RESIDUE, P3-FULLREP-KP, PROOF-IN11}.md
%  and research/v3/species/{C01-CBAR1, C03-BYPASS, X01-L1C, Z06-HOLEA,
%  Z07-HOLEB, Z08-HOLEC, Z09-HOLED, Z10-HOLEE, Z12-FIXC, Y01-KL,
%  Y04-AUDITX01, A02-2PPI2, W95A-2PP, W95C-ORDER, C04-BPRIME}.md
%  House rules honoured here: every substantive claim carries a chip and a
%  file name; no closing symbol is printed anywhere; no Clay claim; [C] and
%  [E] are never written as [T].
% =====================================================================
\begingroup
\renewcommand{\code}[1]{\texttt{\footnotesize\detokenize{#1}}}

\section{The strong-coupling window at $[\mathrm T]$: the residue $I1$, the isolation $I2$, and
the five holes}
\label{ch:uncond-i1}

\noindent The residue theorem of \S\ref{sec:p3} was stated at grade \LT{} in its residue and with
its isolation left open. This chapter is the ledger of what happened to the five holes that the
audit of that theorem attached, and of the two spectral statements the polymer expansion is
supposed to deliver. It states both statements with their windows (\S\ref{sec:u1-theorem}), walks
the five links of $I1$ one line each (\S\ref{sec:u1-five-rings}), closes the five holes in one
table (\S\ref{sec:u1-holes}), settles the accounting constant at its exact value $\bar C\equiv N$
--- one identity refuted, one window re-derived (\S\ref{sec:u1-cbar}), states the four
qualifications and the target that replaces window-extending (\S\ref{sec:u1-d}), instances the
same chain in the $2^{++}$ channel (\S\ref{sec:u1-i2}), and lists what is still owed
(\S\ref{sec:u1-ledger}).

\subsection{The main theorem, and its three side conditions}
\label{sec:u1-theorem}

\begin{claimbox}{cT}
\claimhead{Theorem}{(strong-coupling window: residue and isolation, two channels)} \emph{Let
$\mathcal C_L:=T^3_L\times_a\mathbb Z$ be the cylinder (spatial torus of side $L$, infinite
discrete time, spacing $a$), let $O=\tfrac14\operatorname{tr}F^2$, and let $m_1$ be the lightest
energy with nonzero $O$-overlap. Fix $N$ and let, in the SZZ coupling $\beta_A=2/(g^2N)$ (Wilson: $\beta_W=N^2\beta_A$; at the quoted $N=2$ the two differ by $N^2=2N=4$)}
\begin{equation}
\beta_A\ <\ \beta_\star\ :=\ 3.777\times10^{-3}\ =\ 0.1813\times\tfrac1{48}
\qquad\bigl(\text{i.e. }\beta_W<0.015108=0.725\times\tfrac1{48}\bigr),
\label{eq:u1-beta}
\end{equation}
\emph{(\LC{} on the value: $3.777\times10^{-3}=W(a^\ast)$ is the computed maximum of the
criterion at $a^\ast=0.618$; \eqref{eq:u1-window2} below records the $\hat C=1$, $a=0$ reading
that it replaces. The structure that follows is \LT{}; the window is the one slot of this
theorem whose value is computed, not proved.)}
\emph{Then for every finite $L$, \emph{uniformly} in $L$ --- $L$ is a uniformity axis and
$L\to\infty$ is never taken: the order is fixed-$L$, then $t\to\infty$, then $\inf_L$ ---}
\begin{enumerate}[label=(\roman*)]
\item \emph{(residue, $I1$)} \emph{the lightest $0^{++}$ atom of $O$ satisfies}
      \begin{equation}
      f_{1,L}^2:=\nu_L(\{E_{1,L}\})=\lim_{t\to\infty}e^{E_{1,L}t}C_L(t)\ \ge\ r_0\ >\ 0\,;
      \label{eq:u1-f1}
      \end{equation}
\item \emph{(isolation, $I2$)} \emph{and}
      \begin{equation}
      \Lambda_\star-m_1^2\ \ge\ \delta\ >\ 0,
      \qquad \Lambda_\star:=\inf\operatorname{supp}\bigl(\nu-f_1^2\delta_{m_1^2}\bigr)\,;
      \label{eq:u1-i2}
      \end{equation}
\item \emph{($2^{++}$ channel)} \emph{both (i) and (ii) hold for
      $O_{2^{++}}=\sum_x[2P_{12}-P_{13}-P_{23}]$ with its own $\nu^{(2)}$, $m_2$,
      $f_{1,2}^2$ and $\Lambda_\star^{(2)}$;}
\end{enumerate}
\emph{whence both factors of $\min(\Delta_0,\Delta_2)$ carry a certificate.} \LT{}, \emph{with the
Ursell constant at its exact value $\bar C\equiv N$} (\S\ref{sec:u1-cbar};
\code{F07-ADJCBAR.md}). \emph{The level is fixed by D.1--D.2 of \S\ref{sec:u1-d}: (i)--(iii)
certify the two \emph{factors} of the minimum; the step from them to $m_1>0$ is vacuous (D.1), and
the positivity of the mass inside the window is SZZ/OS's (D.2), not a consequence claimed here.}
\end{claimbox}

\noindent\textbf{The three side conditions, stated in parallel with the theorem and never
dropped.} \emph{(1) Coupling:} $\beta_A<\beta_\star=3.777\times10^{-3}=0.1813\times\tfrac1{48}$, i.e.\ $\beta_W<0.015108=0.725\times\tfrac1{48}$ (\LC{} on the value) --- a
strong-coupling window of a polymer expansion, not a physical coupling.
\emph{(2) Volume:} $L$ finite, but the bounds are
$L$-uniform; no $L\to\infty$ limit is taken, and the infinite-volume atom $f_{1,\infty}^2$ is
\emph{not} claimed (\code{P3-CYLINDER-RESIDUE.md} \S1.4). \emph{(3) Scope:} this is
\textbf{not Clay}. The spacing $a$ is fixed, the continuum limit $a\to0$ is not taken, nothing
crosses to $\beta\approx2.5$, and no statement here mentions the continuum mass gap.
\textbf{Clay remains $0\%$.}

\subsection{Four qualifications (D.1--D.4), and the target that replaces window-extending}
\label{sec:u1-d}

\noindent These four are part of the statement, not commentary on it: each names a step that
\eqref{eq:u1-f1}--\eqref{eq:u1-i2} do \emph{not} license, and each was written as if they did in
some earlier draft.

\begin{enumerate}[label=\textbf{D.\arabic*.},leftmargin=3.4em]
\item \textbf{(\code{G07-KATOADJ.md}) $I2\Rightarrow m_1>0$ is \emph{true but vacuous}.} The side
      condition of $I2$, $\sup\operatorname{spec}\widetilde T<1$, is \emph{verbatim} $m_1>0$
      (\code{B01-RP.md}:12), so the conclusion is written in through the side condition: this is not
      an implication. The \emph{only} legal counterexample needs a \textbf{degenerate vacuum}
      ($r_{\mathrm{iso}}=1-e^{-a\delta}=0.4484>0$ with $m=0$), and that is excluded by the
      Perron--Frobenius uniqueness of the vacuum (\code{G07-KATOADJ.md} \S0 F, \S1.4--1.5). The
      witness of the earlier draft of this chapter, $\nu=\delta_0+\mathbf 1_{[\delta,\delta+1]}$,
      is \textbf{withdrawn}: after the VEV is subtracted $f_1^2=0$, so $I1$ fails; reading
      $\delta_0$ as a connected atom would need a degenerate vacuum; and by the letter of the
      definition the post-vacuum $\sup=e^{-a\delta}$ is the top of an a.c.\ block, so
      $r_{\mathrm{iso}}=0$. It witnesses $m>0\wedge\neg\mathrm{Kato}$, not
      $\mathrm{Kato}\wedge\neg(m>0)$.
      \emph{Scale-free reading:} $r_{\mathrm{iso}}$ and $I2$ are the \textbf{same quantity}
      ($E_2-E_1=\Lambda_\star-m_1^2$; both equal $0.625$ on 2D YM), so the real comparison is
      $r_{\mathrm{iso}}$ (a \emph{distance}) against $m$ (a \emph{position}). What (i)--(iii) do
      certify is the two \emph{factors} of $\min(\Delta_0,\Delta_2)$; the step from that to
      $m_1>0$ is exactly the vacuous step above.
\item \textbf{The $m>0$ credit inside the window is SZZ/OS's, not this chapter's.} For
      $\beta<\beta_1$ the positivity of the mass is Osterwalder--Seiler's cluster expansion, and
      the log-Sobolev radius of Shen--Zhu--Zhu is where $1/48$ comes from; the increment here is
      the residue $I1$ and the isolation $I2$, i.e.\ the removal of the premise ``$0^{++}$ is
      strictly lightest'' (\S\ref{sec:u1-i2}) --- and nothing more.
\item \textbf{$\beta_A<3.777\times10^{-3}$ is a coarse-lattice window, not a continuum statement.} At the
      measured point $a\Lambda(1/48)=0.570$ and $m_1a=18.26$ are pure
      numbers (\code{G05-BETAOFa.md} \S2.2--2.3), so the spacing has a positive lower bound,
      $a\ge18.26/m_1$; the window is a window in the \emph{coupling} at fixed $a$, and no $a\to0$
      is taken anywhere in this chapter. \emph{Value:} $\mathrm{SU}(2)=\mathrm{SU}(3)$ exactly; the
      earlier $0.48/0.41$ split was an artifact of a sign error in the two-loop prefactor
      (\code{G05A.md} \S3.2), \LT{}.
\item \textbf{This is not Clay.} No continuum mass gap is claimed, nothing crosses
      $\beta\approx2.5$, and \textbf{Clay remains $0\%$}.
\end{enumerate}

\noindent\textbf{What the residual gap is, and what it is not.} The residual gap is
$(\mathrm{E\text{-}I2})$: there exist $\delta,c>0$ such that
$\inf_{L,a}(\Lambda_\star-m_1^2)\ge\delta$ and $\inf_L m_1^2\ge c$ --- the quantifier shape that
C5 requires. An earlier draft of this chapter named that gap $\Delta_0>0$; \textbf{that
identification is withdrawn}: $\Delta_0$ lives on the curvature slot
($\delta_{\mathrm{curv}}:=a^2\Delta_0^2+W_{\mathrm{full}}$, dimensionless) while $\Lambda_\star$
lives on the spectral measure $\nu$ (mass squared), and $H1\not\Rightarrow I2$
(\code{T01}, \code{N12-VERDICT.md}); moreover $\Delta_0>0$ fails three of the four model tests
(on 2D YM $\Delta_{\mathrm{dyn}}=-W_{\mathrm{full}}<0$, on M6 a gapless model passes, and the single
plaquette gives $-0.0691$). No \LT{} statement of this chapter is conditioned on it.

\noindent\textbf{The target that replaces window-extending.} A constant floor carried to
$\beta\to\infty$ would give $m_{\mathrm{phys}}\ge r/a\to\infty$ as $a\to0$ --- self-negating
(\code{G05-BETAOFa.md} \S3.2) --- so extending \emph{this} window is impossible in principle, and
the object worth stating is a ratio that survives the continuum limit:
\begin{equation}
\liminf_{a\to0}\ \frac{m_{\mathrm{lat}}(a)}{a\Lambda}\ \ge\ c_0\ >\ 0
\qquad(\Longleftrightarrow\ m/\Lambda\ge c_0),
\label{eq:u1-newtarget}
\end{equation}
with $m_{\mathrm{lat}}$ and $a\Lambda$ in the same scheme and the same normalization. Nothing in
this chapter proves \eqref{eq:u1-newtarget}: everything here is at fixed $a$, and
\eqref{eq:u1-newtarget} is the target the next machine has to reach, not a result.

\begin{v2box}
\textbf{The window, and the one constant that costs it.} The Ursell accounting constant is not an
unknown and not a free bound: it is \emph{exactly} the ratio of the per-tile weight to the
normalized activity,
\begin{equation}
\bar C\ \equiv\ \frac{|w|^{1/n}}{u}\ =\ \frac{c_f}{u}\ =\ N .
\label{eq:u1-cbarN}
\end{equation}
The reason is exact and combinatorial: in a closed tube every link carries one $f$ and one $\bar
f$, so all recoupling multiplicities are $1$ ($\prod\mathrm{mult}=1$) and $|w|^{1/n}=c_f$ exactly,
while $u=c_f/N$; and the animal factor is \emph{already inside} $e\Delta$ ($e$ = labelled trees
per plaquette, $\Delta=20$ = adjacency slots; \code{P3-FULLREP-KP.md}:89 verbatim
$\lambda_a=e\Delta$), so multiplying once more by $\Delta$ double-counts and
$\bar C=\Delta=20$ is not a bound on this slot at all
(\code{F07-ADJCBAR.md} \S2--3; \code{C05-AUDITCBAR.md}). With the interface constant folded
($\hat C=1$), $\bar u=e\Delta\,N\,u$ and the criterion $\bar u<1$ gives
\begin{equation}
\beta_\star=\frac{2N^2}{e\Delta\,N}=0.073576=3.53\times\tfrac1{48},
\qquad \bar u(\tfrac1{48})=0.2831544 ,
\label{eq:u1-window2}
\end{equation}
so $\beta=1/12=0.08333$ lies \emph{outside} the window --- which withdraws every earlier claim
that $\beta=1/12$ was the first provable point (\code{F07-ADJCBAR.md} \S4). \emph{Superseded:}
$\hat C=4$ is forced (\code{M07-WINDOW2.md}), so \eqref{eq:u1-window2} is the $\hat C=1$, $a=0$
reading and not the criterion window; that window is $\beta_A<3.777\times10^{-3}=0.1813\times\tfrac1{48}$,
i.e.\ $\beta_W<0.015108=0.725\times\tfrac1{48}$ (\LC{}, \S\ref{sec:u1-theorem}), so the exclusion
of $1/12$ survives a fortiori. \emph{Window and normalization.} $\hat C=4$ and $\bar C=N=2$ are
uniquely fixed by the KP criterion and the $u$-normalization (\code{M07-WINDOW2.md}), and the
resulting arithmetic is machine-verified in Lean (\code{M02-CBAR.lean}: 25 theorems, zero
\texttt{sorry}, axioms \texttt{[propext, Classical.choice, Quot.sound]}). Hence
$\hat C\bar C=8>7.0633$ and the window is
\[
\beta_A<0.0183940=0.88291\times\tfrac1{48}
\qquad\bigl(\text{equivalently }\beta_W<0.0735759=3.5316\times\tfrac1{48};\ \text{the two differ by }N^2=4\bigr).
\]
The value $1/48=0.0208333$ therefore lies \emph{outside} the window, and $\beta=1/12$ lies far
outside.

\smallskip
\noindent\textbf{Correction to the window's number (its number, not $I1$'s content).} The radius
displayed above is \emph{not} the KP criterion evaluated. $0.018394=1/(e\Delta)$ is the $a=0$
activation root, written $\lambda_a\Phi(0)=1$ and evaluated at leading order $\Phi(0)\simeq\beta$
(keeping $\Phi(0)=\beta+\tfrac38\beta^2$ moves the root to $0.0182686$, the $0.018269$ of
\code{UC-COMPLETE.md:122}) --- a radius of convergence, and at $a=0$ the criterion
$S_x(a)\le a$, rewritten $\lambda\Phi(a)\le a/(1+a)$, has right-hand side $0$ while the left-hand
side stays $\ge|w|>0$, so that value is \textbf{illegal} as a criterion window. The criterion at
$a>0$ is read off $W(a)=a/[(1+a)\lambda\varphi(a)]$, $\varphi(a):=\Phi(a,\beta)/\beta$: for a pure
exponential envelope its maximum is the closed-form $a^\ast=(\sqrt5-1)/2=0.618$, and the measured
$\Phi$ (ratio $1:2.7508:7.7238$, slightly above $1:e:e^2$) moves the numerical maximum to
$0.6150$ with the same value to four digits, so the window is
\[
\beta_A\ <\ W(a^\ast)\ =\ 3.777\times10^{-3}
\qquad\bigl(\beta_W<1.511\times10^{-2}\ \text{by }N^2=4\bigr),
\]
with $\Phi(0)=0.0185211$, $\Phi(1)=0.0509474$, $\Phi(2)=0.1430532$ and $\lambda=e\Delta=54.366$
(\code{ADJ3.md}, \code{ADJ4.md}; \code{w7/adj4\_kp\_a.log} \S2--3). The number printed here thus
\textbf{over-reports by} $0.018269/0.003777=4.84\times$ --- by $0.018394/0.003777=4.87\times$ for
this printing's own $0.018394$. Every ``outside the window'' claim below survives, $1/48$ and
$1/12$ lying outside both radii. What does \emph{not} change is the \emph{content} of $I1$: the five
rings, the five holes A--E closed, $\bar C\equiv N$, $\hat C=4$, $[KL]$, the 25 Lean theorems. Only
the window's \emph{value} drops to \LC{} (the maximum at $a^\ast$ is computed --- a grid plus a fine
scan on a representation sum truncated at $n\le300$ --- not proved), while the structure of $I1$
stays \LT{}.

\emph{Historical
note.} Earlier drafts carried a whole table of bounds on this slot --- an exponential bound in
$\Delta u$, a bound by $e$, a bound by the adjacency degree $\Delta$, and finally a direct bound
from which $\bar C$ had supposedly disappeared. The table collapses into the single line
\eqref{eq:u1-cbarN}: every entry is \textbf{withdrawn}, each of them having lain \emph{below} the
exact value $N$, and the $\bar C$-free route is withdrawn at its input --- $|w_p|\le u$ is off by
the factor $N$ (\code{F07-ADJCBAR.md} \S1.4) --- so the wider window it produced goes with it.
\end{v2box}

\noindent\textbf{Two conventions fixed here, once.} \emph{(a) Units.} $\nu_L=\sum_nf_{n,L}^2
\delta_{E_{n,L}}$ is a measure in the \emph{energy} variable, so its atom sits at
$E_{1,L}=m_{1,L}$; \S\ref{sec:p3} writes the second support point as a mass squared, and the two
readings are related by $\Lambda_\star-m_1^2=2m_1\delta_A+\delta_A^2$, so a positive window in
one variable is a positive window in the other (\code{X01-L1C.md} \S3 row~8, after
\code{X02-THEOREMA.md} \S4). \emph{(b) Normalization of the residue.} With hole~D's bookkeeping
the quantity carrying an $L$-independent $r_0$ is the density $\bar f_1^2:=f_{1,L}^2/V_s$,
$V_s=L^3$; the bare atom weight is extensive,
$f_{1,L}^2=V_s\,a_0(1-O(u))$ (\code{Z09-HOLED.md} \S3). Statement \eqref{eq:u1-f1} is written in
the compact form $f_{1,L}^2\ge r_0$ with the standing understanding that either $r_0$ carries
$V_s$ or $f_1^2$ is read normalized; the two are one inequality multiplied by $V_s$, and the only
form that is \emph{not} available is a bare extensive $f_{1,L}^2$ against an $L$-independent
$r_0$ (\code{Z09-HOLED.md} \S5). The constants are $\tfrac12(2N)^{-2}$ at each $L$ and
$\tfrac14(2N)^{-2}$ after the $\inf_L$, both $\beta$-independent (\code{Z10-HOLEE.md} \S5).

\subsection{$I1$ in five rings: four walls, KP, resummation, spectral bridge, $[KL]$}
\label{sec:u1-five-rings}

The proof of \eqref{eq:u1-f1} is a chain of five links. Each link is named with the file that
contains it; each skeleton is one line, because the work is in the file and not in this page.

\begin{enumerate}[label=\textbf{R\arabic*.},leftmargin=2.6em]
\item \textbf{Four walls} (\code{P3-FOURWALL-DERIVATION.md} \S1.1--1.5; Gupta (14.20)).
      \emph{Skeleton:} each of the source plaquette's four boundary links carries an open
      fundamental index pair and $\int dU\,U_{ij}=0$ forbids covering it alone; two unit
      plaquettes share at most one link, so there are $\ge4$ independent strands; only temporal
      plaquettes advance time, so each strand crosses $t/a$ slices at $\ge1$ plaquette each;
      hence $|P|\ge4t/a$, with equality exactly for the straight four-walled tube, whose weight
      is $a_0u^{4t/a}=a_0e^{-m_1t}$, $m_1=w\ln(1/u)/a$, $w=4$. \LT{}
\item \textbf{KP in the full representation} (\code{P3-FULLREP-KP.md} \S2.1--2.3, \S4.1).
      \emph{Skeleton:} Haar is a probability measure, so $|\chi_r(U)|\le\dim r$; the multiplicity
      cap $c_r\le\beta^{|r|}/(|r|!\,\dim r)$ cancels \emph{exactly} the $\dim r$ that the
      contraction would cost, leaving $c_r\dim r\le\beta^{|r|}e^{\beta}/|r|!$, whose factorial
      beats the sub-exponential rep-label entropy; hence
      $\sum_{P\ni x}|w(P)|e^{a|P|}\le a$ whenever $uK<1$, with $K\le6.1\times10^{2}N$ and
      plaquette adjacency $\Delta=4(6-1)=20$ (not $12$). \LT{}
\item \textbf{Resummation of the decorations} (\code{X01-L1C.md} \S1--2,
      \code{Y04-AUDITX01.md} \S1--3; re-coupling and activity fixed in
      \code{Z08-HOLEC.md}, \code{Z12-FIXC.md}). \emph{Skeleton:} group clusters by
      \emph{attachment span} $\alpha$ (not by support); the tube is \emph{open}, so the number of
      placements is the interval count $(T-\alpha+1)_+$; time translation gives $A(\alpha)$ and
      the exact identity
      $\log Z_{\mathrm{dec}}=\sum_\alpha(T-\alpha+1)_+A(\alpha)
      =\varphi_{\mathrm{dec}}T+\log D_\infty+\varepsilon_T$ with
      $\varepsilon_T=\sum_{\alpha>T}(\alpha-1-T)A(\alpha)$; the shift
      $m:=\alpha-1-T$ moves weight and exponent together, so the tail is
      $|\varepsilon_T|\le B\bar u^{T}$ with $B=(\bar u/(1-\bar u))^{3}$ and \emph{no} $T$
      prefactor; the activity is the interface-corrected $w_{B_T}=w\prod_{\ell}R(\ell)$, whose
      envelope is $\hat C=4$ (hole~C). \LT{}
\item \textbf{Spectral bridge} (\code{P3-CYLINDER-RESIDUE.md} \S1.2--1.4, \S2.1(b), \S2.4).
      \emph{Skeleton:} at finite $L$ the transfer matrix has a continuous kernel on a compact
      configuration space, hence is Hilbert--Schmidt, hence compact, hence has purely discrete
      spectrum --- \emph{every} spectral point is an atom; and
      $f_1^2=\lim_te^{E_{1,L}t}C_L(t)$ uses only $\delta_L=E_{2,L}-E_{1,L}>0$ \emph{at that $L$}.
      The residue is a \emph{coefficient} and the isolation is an \emph{eigenvalue spacing}; this
      is why $m_1$ only \emph{locates} the atom that is being read and enters no inequality of
      the chain. \LT{}
\item \textbf{$[KL]$: positivity of the measure} (\code{P3-CYLINDER-RESIDUE.md}:67--69;
      \code{Y01-KL.md} \S4.1). \emph{Skeleton:} $\nu_L\ge0$ is Osterwalder--Seiler reflection
      positivity and $f_{n,L}^2\ge0$ is the spectral theorem, so the K\"all\'en--Lehmann
      representation on the cylinder is a \emph{theorem} --- Osterwalder--Seiler (1978) plus
      finite-$L$ compactness plus the spectral theorem, with \textbf{no Yang--Mills dynamics in
      it} --- and not an assumption; since the statement quantifies over the Hermitian,
      gauge-invariant, time-translation-covariant $O$, it transfers verbatim to $O_{2^{++}}$.
      \LT{}
\end{enumerate}

\noindent\textbf{One residual, named rather than absorbed.} The only place where a physical input
could still enter is the identification of the leading polymer concentration with the leading
atom --- i.e.\ that the two exact representations of $C_L(t)$ are the \emph{same} function
(\code{P3-CYLINDER-RESIDUE.md} \S2.1(d)). The $t$-uniform dressed-mass/amplitude resummation
that this needs is what \code{P3-DRESSED-MASS.md} \S1--4 supplies, and hole~C is its last piece;
\code{Y01-KL.md} \S4.3 therefore keeps one residual $[\mathrm{R}]$ on the \emph{verbatim} equality
$a_0D_\infty=|\langle\Omega|O|\mathbb{G}\rangle|^2$. The floor \eqref{eq:u1-f1} does not need
that equality, and neither does the $I2$ chain of \S\ref{sec:u1-i2}; it is the one debt this
chapter refuses to pay by relabelling.

\noindent\textbf{And one second residual, in the same spirit.} The floor is \LT{} with the Ursell
constant at its exact value $\bar C\equiv N\equiv2$, not with "$\bar C$ removed": the summation
identity that was supposed to give $\bar C=1$ is \emph{false} (\S\ref{sec:u1-cbar},
\eqref{eq:u1-false-id}), and what replaces it is an exact identity $\bar C=c_f/u=N$
(\code{F07-ADJCBAR.md} \S3). A constant whose value is proved is a different object from a
constant that is assumed to be $1$, and the labels of this chapter keep them apart. The tree
factor $\Delta=20$ and the weight ratio $\bar C=N$ live in different slots; the per-cluster value
$2.2209$ quoted below belongs to the first slot and is \emph{not} a lower bound for the second.

\subsection{The five holes and how each was closed}
\label{sec:u1-holes}

\begin{center}
\footnotesize
\setlength{\tabcolsep}{4pt}
\renewcommand{\arraystretch}{1.22}
\begin{tabularx}{\textwidth}{@{}l>{\raggedright\arraybackslash}p{3.15cm}%
  >{\raggedright\arraybackslash}X>{\raggedright\arraybackslash}p{3.4cm}@{}}
\toprule
\textbf{Hole} & \textbf{What it was} & \textbf{How it is closed} & \textbf{Source} \\
\midrule
A & $u=c_f$ vs.\ $c_f/N$: a discrepancy of $N^{4T}$ in the tube weight
  & \textbf{$u:=c_f/N$}, i.e.\ one $1/N$ per tile, verbatim from Gupta (14.18)/(14.21)
    ($\beta/2N$ from the expansion $\times$ $1/N$ from the integration $=\beta/2N^2$); and
    \textbf{$A^{\mathrm{dressed}}$ \emph{is} $f_1^2$}, so there was never a second quantity
  & \code{Z06-HOLEA.md} \S1--5 \LT{} \\
B & whether $\delta\ge\kappa/3$ lands inside $I1$
  & \textbf{out}: neither $I1$ nor $D_\infty$ contains $\delta$, and $D_\infty$ is not a
    function of $\delta$ --- writing $\delta$ into the definition of $D_\infty$ already fails at
    $T=0$; the reachable rate is $\delta=\ln(1/\bar u)/a=\kappa/4-\ln(\hat Ce\Delta\bar C)/a$,
    strictly below $\kappa/4$
  & \code{Z07-HOLEB.md} \S2--4; \code{X01-L1C.md} \S2.1 \LT{} \\
C & $Z_{\mathrm{dec}}$ re-coupling is not block-multiplicative:
    $\mathrm{Hom}(1,f\bar f f\bar f)=2\ne1^2=1$; cap-type decorations have $w(\gamma)=0$ while
    $w(B_T\cup\gamma)\ne0$
  & \textbf{$\hat C=2$ exact at each factorization point, envelope $\hat C=4$} per plaquette;
    then $\delta=\kappa/4-\ln(\hat Ce\Delta\bar C)/a$ and, at $\beta=1/48$ with $\bar C=N=2$,
    $\delta\,a=1.2618>0$ (interface constant folded, $\hat C=1$) or $0.5686>0$ ($\hat C=2$)
    --- the $0.56862$ of the earlier draft was the $\bar C=1$ value and is
    superseded; the interface re-coupling is absorbed into the
    \emph{activity} and, through it, into $\varphi_{\mathrm{dec}}$, never into $\varepsilon_T$
  & \code{Z08-HOLEC.md} \S2--4; \code{Z12-FIXC.md} \S1--3; \code{F07-ADJCBAR.md} \S4 \LT{} \\
D & summing $O$ over the whole slice makes $f_1^2\propto L^3$, so a bare $f_1^2$ cannot be
    compared with an $L$-independent $r_0$
  & \textbf{$r_0$ carries $V_s$} (or $f_1^2$ is read normalized); and $a_0$'s $L$-independence is
    \emph{already proved} --- positive geometric sum over transverse positions, wrap corrections
    $e^{-cL}$, no entropy $\mu^T$
  & \code{Z09-HOLED.md} \S3--5 \LT{} \\
E & $r_0=\tfrac14(2N)^{-2}$ vs.\ $\tfrac12c_1u^{p}$: two $\beta$-scalings for one constant
  & \textbf{$p$ is the operator-coefficient exponent and $p=0$}; $u$ lives only in the correction
    factor, $r_0=\tfrac12a_0(1-O(u))$, $\beta$-independent. At $\beta<1/48$ the $u^{p}$ form with
    $p\ge1$ is $2$--$6$ orders of magnitude too small and misstates $\beta\downarrow0$
  & \code{Z10-HOLEE.md} \S3--5 \LT{} \\
\bottomrule
\end{tabularx}
\end{center}

\noindent Two audit corrections inside the reconstruction, recorded because a witness that does
not witness is worse than none: (i) hole~C's original cap witness (a decoration giving a
\emph{single} $f$ at the shared link) is zero on \emph{both} sides, since
$\mathrm{Hom}(1,f\bar ff)=0$; the genuine witness is a decoration giving
$\mathrm{adj}$, where $w(\gamma)=0$ but $w(B_T\cup\gamma)=1$ (\code{Z12-FIXC.md} \S1.3). (ii)
The cap witness of \code{Z08-HOLEC.md} §2.2 used the wrong arithmetic
($\dim\mathrm{Hom}(f,f)=1$); the qualitative failure survives, the number does not.

\subsection{The accounting constant: $\bar C\equiv N$ (one refuted identity, one window)}
\label{sec:u1-cbar}

\noindent Everything in \S\ref{sec:u1-five-rings} that is not the KP estimate itself is a
statement about \emph{counts}, and every count carried an Ursell/tree constant $\bar C$ --- which
\code{A01-CBAR.md} could bound only by $\bar C\le\Delta=20$, shrinking the window by $11.3\times$.
An earlier draft of this ledger claimed $\bar C=1$ at the summation level and declared the
constant removed. \textbf{That claim is withdrawn in this printing}: the identity it rested on is
false, and the honest grade is \LT{} with the Ursell constant at its exact value $\bar C\equiv N$.
What survives is enough for the window, and the bookkeeping is set out below in the order the
audit settled it.

\subsubsection*{The identity that was hoped for, and the exact counterexample}

Let $G_\Gamma$ be the plaquette-adjacency graph of a connected cluster $\Gamma$. The Ursell
(Mayer) coefficient is an alternating sum over connected edge subsets,
\begin{equation}
\varphi^U(\Gamma)=\sum_{A\subseteq E(G_\Gamma),\,(\Gamma,A)\ \text{connected}}(-1)^{|A|},
\label{eq:u1-ursell}
\end{equation}
so its natural control is a \emph{tree} count: the Penrose tree-graph inequality
$|\varphi^U(\Gamma)|\le\tau(G_\Gamma)$ (the number of spanning trees) is \LT{} --- a four-line
delete--contract induction, plus $0$ violations on $875$ enumerated graphs
(\code{C07-AUDIT.md} \S4.2--4.3). The step that was \emph{hoped} for is
\begin{equation}
\boxed{\ \sum_{\Gamma\ni x,\ \Gamma\ \text{connected}}\tau(G_\Gamma)\prod_{p\in\Gamma}a_p
\;=\;\sum_{T\ni x,\ T\ \text{a tree}}\ \prod_{p\in V(T)}a_p\ }\qquad\text{--- \textbf{false}.}
\label{eq:u1-false-id}
\end{equation}
The smallest witness is the triangle: with $V=\{x,y,z\}$ and edge activities $a_0,a_1,a_2$,
\[
\text{LHS}=a_0a_1+a_0a_2+a_1a_2+3a_0a_1a_2,\qquad
\text{RHS}=a_0a_1+a_0a_2+a_1a_2,
\]
so $\text{LHS}-\text{RHS}=3a_0a_1a_2\ne0$: the triangle's $\tau=3$ spanning trees are counted
three times on the left, and the right side has no cubic term at all. At $n=4$ ($K_4$) the
discrepancy is a $22$-term polynomial of degrees $[4,6]$, with the tree terms cancelling exactly.
\textbf{The tree sum does not absorb the animal factor; the animal factor comes back}, and it
comes back as the product over non-tree edges (\code{C07-AUDIT.md} \S1.1--1.2, \code{[T]}: the
counterexample is an exact identity of integers).

\subsubsection*{The corrected identity, and the exact value of $\bar C$}

Grouping by spanning tree gives the correct statement,
\begin{equation}
\sum_{\Gamma\ni x}\tau(G_\Gamma)a^{\Gamma}
\;=\;\sum_{T\ni x}a^{T}\prod_{e\notin T}(1+a_e),
\label{eq:u1-tree-id}
\end{equation}
verified symbolically at $n=3,4$ and in integers at $n=4,5$ (unit activities: ratios
$2,8,64=2^{\#\text{non-tree edges}}$) (\code{C07-AUDIT.md} \S1.2, \LT{}). There is also a sharper
structural reading: $\varphi^U$ factorizes over $2$-connected blocks,
$\varphi^U=(-1)^{\#\text{bridges}}\prod_B\mu(B)$, and factorial growth occurs only on cliques,
while a link carries exactly $2(d-1)=6$ plaquettes and any two of them share that link --- so
$K_6$ is the largest clique and $(n-1)!$-growth stops at $n=6$
(\code{C02-URSELL.md} \S1, \LT{}). What the corrected identity does \emph{not} do is make
$\bar C$ small: the non-tree product is a genuine animal factor, and it has to be paid. The
resolution the audit finally reached is that $\bar C$ is not a bound to be optimised at all but a
\emph{ratio}, and that its value is exact:
\begin{equation}
\bar C\ \equiv\ \frac{|w|^{1/n}}{u}\ =\ \frac{c_f}{u}\ =\ N
\qquad(N=2:\ \bar C=2),
\label{eq:u1-cbar-exact}
\end{equation}
because in a closed tube each link carries one $f$ and one $\bar f$, so every recoupling
multiplicity is $1$ ($\prod\mathrm{mult}=1$, verified on the cube and on the four-wall tube at
$n=6,4,8,12$) and $|w|^{1/n}=c_f$ exactly, while the normalized activity is $u=c_f/N$. The
consistent sharpened account is $\Phi(0)/u=2N^2$ (with $\Phi(0)=\sum_{r\ne0}c_rd_r=\beta=2N^2u$,
the $f$ and $\bar f$ branches), so $\bar C\in[N,2N^2]$
(\code{F07-ADJCBAR.md} \S1.4 and \S3, \LT{}; \code{C05-AUDITCBAR.md}).

\begin{qbox}
\textbf{What is \emph{not} available, recorded so that it is not tried again.} Two earlier
readings are withdrawn. (i) $\bar C=1$ at the summation level, by the refuted identity
\eqref{eq:u1-false-id}. (ii) Every bound of the form $\bar C\le(\text{const})$ that earlier drafts
tried on this slot: each of them lay \emph{below} the exact value $N$ and is therefore not a bound
on the slot at all --- the slot is closed by \eqref{eq:u1-cbar-exact}, not by a bound. Two
things do survive and are worth keeping. The tree-factor slot is genuinely $N$-free:
$\Delta=20$ is a four-dimensional adjacency degree, and the animal count $e\Delta$ already
contains the labelled-tree factor $e$, so the per-cluster supremum $|\varphi^U|^{1/n}=2.220906$
(attained at $K_6$, with $120=(6-1)!$) is a statement about $\tau$\emph{-counting}, not about
$\bar C=c_f/u$; the two share a symbol and nothing else (\code{F07-ADJCBAR.md} \S3, \S7). And the
$\tau$ route cannot be pushed to $3.5$ by any refinement: the $L=2$ periodic box cluster has
$\tau=3.3559\times10^{116}$, i.e.\ $\tau^{1/96}=16.361$, so $\Delta=20$ is nearly saturated and
the looseness is in the object $\tau$, not in the algebra (\code{C02-URSELL.md} \S2.2, \LT{}).
\end{qbox}

\subsubsection*{Consequence: the window, re-derived}

With $\bar u:=\hat C\,e\Delta\,\bar C\,u$ and $\bar C=N$ exact, the lossless condition
$\hat C\bar C<2N^2\cdot48/e\Delta=7.063285$ ($N=2$) becomes a condition on $\hat C$ alone:
$\hat C N<7.0633\iff\hat C<3.5316$, and the corresponding radius is
\begin{equation}
\beta_\star=\frac{2N^2}{e\Delta\,N}=0.073576=3.53\times\tfrac1{48}
\qquad(\hat C=1;\ \text{superseded by }\hat C=4\text{; see the note below}),
\label{eq:u1-window}
\end{equation}
with $\bar u(1/48)=0.2831544$, $B=(\bar u/(1-\bar u))^3=0.0616$ and $\delta\,a=\ln(1/\bar u)
=1.2618$ at $\beta=1/48$; $\kappa/4\cdot a=4.564348$ is untouched, being $\bar C$-free. Two
readings of the interface constant are worth recording, because they differ in kind: with the
\emph{exact local} $\hat C=2$ of hole~C the radius is $0.036788=1.766\times\tfrac1{48}$ (still
containing $1/48$), while with the crude per-plaquette envelope $\hat C=4$ it is
$0.018394=0.883\times\tfrac1{48}$, which no longer contains $1/48$ and at which $\bar u=1.1326>1$
already at $\beta=1/48$. The window quoted by the theorem is the criterion value $W(a^\ast)$ of the next paragraph, not the $\hat C=1$ accounting of
\code{F07-ADJCBAR.md} \S4; the $\hat C=4$ reading is carried in \S\ref{sec:u1-ledger} as a debt
rather than hidden (\LT{}+\LE{}). \emph{Correction:} both radii in this paragraph are $a=0$ roots
($\lambda_a\Phi(0)=1$, i.e.\ convergence radii), not values of the criterion; the criterion at
$a^\ast=0.618$ gives $3.777\times10^{-3}$ (\S\ref{sec:u1-theorem}, \LC{}), and what this paragraph
keeps is the pair $\hat C=4$, $\bar u(1/48)=1.1326>1$.

\noindent\textbf{Note after \code{M07-WINDOW2.md}.} The $\hat C=1$ reading of the condition displayed
above is superseded: $\hat C=4$ is forced by the KP criterion together with the $u$-normalization,
so $\hat C\bar C=8>7.0633$ and the window is the one quoted in \S\ref{sec:u1-theorem}, namely
$\beta_A<0.88291\times\tfrac1{48}$, i.e.\ $\beta_W<3.5316\times\tfrac1{48}$. The number
$0.073576$ survives only as the $\hat C=1$, $a=0$ value of that pair in $\beta_W$. Both radii
are $a=0$ quantities; the criterion window at $a^\ast$ is $\beta_A<3.777\times10^{-3}$, i.e.\
$\beta_W<0.015108=0.725\times\tfrac1{48}$ (\S\ref{sec:u1-theorem}, \LC{}).

\subsubsection*{The route that would have removed $\bar C$ outright --- withdrawn}

Earlier drafts of this ledger quoted a direct bound on the anchored decoration sum as an
alternative route that avoids $\bar C$ altogether:
\begin{equation}
A(\alpha)=\sum_{\Gamma:\ \alpha(\Gamma)=\alpha,\ \mathrm{att}\,\Gamma=[0,\alpha]}
\varphi^U(\Gamma),
\qquad
|A(\alpha)|\le N_{\mathrm{piv}}\frac{(e\Delta u)^{\alpha}}{1-e\Delta u},
\label{eq:u1-Aalpha}
\end{equation}
with $N_{\mathrm{piv}}\le4\times6=24$ (\code{C03-BYPASS.md} \S2.4). \textbf{That route is
withdrawn.} Its input $|w_p|\le u$ is off by exactly the factor \eqref{eq:u1-cbar-exact}: the
per-plaquette weight is $c_f=Nu$, so the correct bound reads
$N_{\mathrm{piv}}(e\Delta Nu)^{\alpha}/(1-e\Delta Nu)$, whose ratio is
$e\Delta Nu=1.1326>1$ already at $\beta=1/12$ — the criterion fails there, and the window it had
claimed is gone (\code{F07-ADJCBAR.md} \S1.4 and \S4, \LT{}). The $\bar C$-free slot was an
illusion produced by moving $\bar C$ from the tree factor into the weight ratio without noticing.
Nothing in this chapter uses \eqref{eq:u1-Aalpha}; what it asserts is \eqref{eq:u1-window}.
\emph{Correction:} $e\Delta Nu=1.1326$ at $\beta=1/48$ is the $a=0$ reading, so the failure asserted
here is a fortiori; the criterion's own window is $3.777\times10^{-3}$
(\S\ref{sec:u1-theorem}, \LC{}).

\subsection{$I2$, and the two lightest channels}
\label{sec:u1-i2}

\noindent $I2$ is \emph{not} an eigenvalue statement of the polymer expansion --- that expansion
controls a weight (a coefficient), not a spacing, which is exactly why
\code{PROOF-P3-UNCONDITIONAL.md} \S5 writes ``\textbf{NOT obtained:} isolation $I2$''. The
isolation comes from a different object, the two-exponential tail of $C_L(t)$.

\begin{claimbox}{cC}
\claimhead{Theorem}{($I2$ via the two-exponential chain, one instantiation per channel)}
\emph{On $\mathcal C_L$, for $\beta_A<\beta_\star$ and every finite $L$, the chain}
\[
(L1)\ \xrightarrow{\ \mathrm{[T]}\ }\ \text{algebra}\ \xrightarrow{\ \mathrm{[T]}\ }\
\text{Theorem~A}^{\Leftarrow}\ \xrightarrow{\ \mathrm{[T]}\ }\ I2
\]
\emph{closes, where each arrow is a \LT{} link: $(L1)$ is $|\log D_{\mathrm{fin}}(T)-\log
D_\infty|\le B\bar u^{T}$ with $\bar u=\hat C e\Delta\bar C u<1$, the algebra is
$D_{\mathrm{fin}}/D_\infty-1=e^{\varepsilon}-1$ with $|e^{\varepsilon}-1|\le2|\varepsilon|$,
and Theorem~A$^{\Leftarrow}$ is the measure-theoretic equivalence
$\text{atom}(I1)\wedge I2\iff$ sharp two-exponential decay
(\S\ref{sec:reslevel}). Hence $\Lambda_\star-m_1^2=2m_1\delta+\delta^2>0$.} \LT{}
\end{claimbox}

\noindent Three things about this chain are load-bearing and are stated separately.

\emph{(a) Positivity transfers because $O$ is universally quantified.} $\nu^{(2)}_L\ge0$ is the
same $[KL]$ as in ring 5: the statement is proved for \emph{any} Hermitian, gauge-invariant,
time-translation-covariant $O$, and $O_{2^{++}}$ is one. $\nu\mapsto\nu^{(2)}$ is therefore a
free rewrite at the level of \emph{definitions} (\code{A02-2PPI2.md} \S1.2); the resulting
proposition $I2^{(2)}$ is logically independent of $I2^{(0)}$ but needs no separate proof,
because it is the same machine instantiated twice.

\emph{(b) The parameters coincide verbatim, which is stronger than leading degeneracy.} The five
rings eat only $w=4$ and the per-tile weight $u$; neither knows the channel. Hence
\begin{equation}
\kappa^{(2)}=4\ln(1/u)/a=\kappa^{(0)},\qquad
\delta^{(2)}=\delta^{(0)}=\frac{\ln(1/\bar u)}{a}
=\frac{\kappa}{4}-\frac{\ln(\hat C\,e\Delta\,N)}{a},
\label{eq:u1-kappa}
\end{equation}
with $\bar u,B$ likewise equal. This is \emph{not} the statement $m_2/m_0\to1$: $\kappa$ is the
bare tube rate and $m^{dr}=\kappa-f_{\mathrm{dec}}$ is the mass, and the coincidence is at the
level of the rate \emph{parameter}, a combinatorial fact about the operator's per-slice plaquette
content (\code{Y04-AUDITX01.md} \S4.1: the $4$ is the source plaquette's four edges, not $d=4$
and not a face count). \LT{}

\emph{(c) The chain needs $m_2>0$, not $m_2=m_0$.} Theorem~A is a two-parameter statement: each
channel uses its own mass in $e^{-mt}$ and its own threshold in $\Lambda_\star-m^2$. What the
chain needs is $m_2>0$, which is Osterwalder--Seiler cluster expansion per channel at
$\beta<\beta_0$ (\code{W95A-2PP.md} \S4, \code{[T]}), not the equality of the two masses --- and
that equality is in any case only a leading-order statement whose $O(1)$ subleading splitting is
neither computed nor sign-fixed (\code{PROOF-IN11.md}, the section ``Why the core ordering fails
even here''; \code{W95C-ORDER.md} \S2). So the requirement is \emph{weaker} than degeneracy, and
the degeneracy's failure could not break the chain.

\begin{claimbox}{cT}
\claimhead{Corollary}{(both lightest factors certified, and membership)}
\emph{With \eqref{eq:u1-f1} instantiated in both channels and \eqref{eq:u1-i2} following from the
chain above, both factors of $\min(\Delta_0,\Delta_2)$ carry a certificate, so $m>0$ in the
window no longer rests on the non-theorem "$0^{++}$ is strictly lightest". The membership
premise needed to exclude the remaining channels is $(b')_{\mathrm{IN-11}}$: no bulk state below
$\delta$ lies outside $\{0^{++},2^{++}\}$. On $\beta<\beta_1$ it is a theorem ---
\code{PROOF-IN11.md}:19, verbatim ``This proves $(b')$'' --- obtained from the convergent OS/SZZ
expansion plus a finite cubic-group lemma ($a_c\ge2$ for every other $J^{PC}$, hence
$m_c-m_{0^{++}}\ge(a_c-1)\ln(1/u)\to+\infty$). Beyond $\beta_1$ it is physical input and is
\emph{not} claimed (\code{C04-BPRIME.md} \S5).} \LT{}, \emph{with the exact Ursell constant
$\bar C\equiv N$} in the $I1$ factor (\S\ref{sec:u1-cbar}); the $I2$ factor carries the same
constant explicitly through $\bar u=\hat C e\Delta N u$.
\end{claimbox}

\noindent The four limits of this corollary, none of them optional: it is \emph{not} the first
proof of $m>0$ at strong coupling, which is Osterwalder--Seiler; its increment is that the two
factors of the minimum are certified, i.e.\ that the premise "$0^{++}$ strictly lightest" is no
longer used. It does not cross $\beta\approx2.5$; it takes no $a\to0$; and it keeps the
normalization of \S\ref{sec:u1-theorem}(b) and the membership $(b')_{\mathrm{IN-11}}$
(\code{W95C-ORDER.md} \S4, \code{A02-2PPI2.md} \S4(a)).

\subsection{The remaining technical debts}
\label{sec:u1-ledger}

\noindent These are recorded as numbers, not as prose, because each of them has bitten a previous
draft.

\begin{itemize}[leftmargin=1.5em]
\item \textbf{The $\bar C$ ledger, in one line.} $\bar C\equiv|w|^{1/n}/u=c_f/u=N$ \emph{exactly}
      --- a proved value, not a bound, and not an unknown (\code{F07-ADJCBAR.md} \S1.4, \S3).
      Consequences, all of them load-bearing: the summation identity that gave $\bar C=1$ is
      false (\code{C07-AUDIT.md} \S1.1, witness $3a_0a_1a_2$); the later bounds tried on this slot
      are withdrawn rather than merely superseded, each having lain \emph{below} $N$; and the
      window is
      $\beta_A<W(a^\ast)=3.777\times10^{-3}=0.1813\times\tfrac1{48}$, i.e.\
      $\beta_W<0.015108=0.725\times\tfrac1{48}$ (\LC{}, \S\ref{sec:u1-theorem}), so
      $\beta_W=1/12=0.08333$ is $5.5\times$ \emph{outside} it. \emph{Debt, stated rather than
      hidden:} the two remaining radii are $a=0$ roots --- the exact local $\hat C=2$ of hole~C
      gives $1.766\times\tfrac1{48}$ (still containing $1/48$), and the crude per-plaquette
      envelope $\hat C=4$ gives $0.883\times\tfrac1{48}$, which does not. (Those two are $a=0$
      roots; the criterion window is $W(a^\ast)=3.777\times10^{-3}$, \S\ref{sec:u1-theorem},
      \LC{}.)
\item \textbf{The $\bar C$-free route is withdrawn.} The direct bound on $A(\alpha)$ rested on
      $|w_p|\le u$, which is off by the factor $N$; with the correct weight its ratio is
      $e\Delta Nu=1.1326>1$ at $\beta=1/12$ and the claimed window is gone
      (\code{F07-ADJCBAR.md} \S1.4, \S4). Withdrawn (\code{F07-ADJCBAR.md}), not outstanding.
\item \textbf{$\beta_1$ scales like $cN$, not $c/N$.} With $u=c_f/N$ and $K\le591e^{\beta}N$ the
      criterion $uK<1$ reads $\beta e^{\beta}<N/295.6$, i.e.\
      $\beta_1(N)=W(N/295.6)$ and, at small argument, $\beta_1=cN$ with
      $c=1/295.6=3.383\times10^{-3}$. The direction matters: $uK\propto\beta/N$, so a larger $N$
      gives a \emph{larger} radius, and $\min(1/48,cN)$ is taken by $1/48$ for $N\le6.3$ --- at
      $N=2,3$ this KP radius does not enter the minimum at all. The physical window is
      $N$-independent after all: $g^2_1=2N/\beta_1\approx594+O(\ln N)$
      (\code{LOOP-RUNBOOK.md} \S13.48, \code{V06-RADIUS.md}; \LT{} algebra, \code{[E]} constants).
\item \textbf{$1/48$ is a method radius, not a spectral wall.} It is Shen--Zhu--Zhu's
      log-Sobolev radius $|\beta|<1/(16(d-1))=1/48$; no statement about the mass gap follows from
      it, and no wall is located by it. What it does is bound the window in which the convergent
      expansion is available (\code{LOOP-RUNBOOK.md} \S13.48).
\item \textbf{$\bar C$ is $N$-independent.} $\Delta=20$ is a four-dimensional plaquette degree and
      $e\Delta=54.3656$ carries no $N$; the $N$-dependence of the radius sits entirely in
      $K\le6.1\times10^{2}N$ (\code{P3-FULLREP-KP.md} \S4.1--4.2). So "the constant got worse"
      and "the window shrank with $N$" are different statements and only the second is false.
\item \textbf{At $N=2$ the character coefficient doubles.} Because $f\cong\bar f$, the two
      orientations add, so $c_f=\beta/N$ rather than the generic $\beta/(2N)$, and
      $u=c_f/N=\beta/N^2$ --- Gupta's own footnote to (14.18). Every number in this chapter is
      quoted at $u=\beta/2N^2$ with the $N=2$ caveat applied where it applies
      (\code{Z06-HOLEA.md} \S1, \S3).
\item \textbf{One residual $[\mathrm{R}]$, already named.} The verbatim equality
      $a_0D_\infty=|\langle\Omega|O|\mathbb{G}\rangle|^2$ (\code{Y01-KL.md} \S4.3(3)); the floor
      and the $I2$ chain do not use it.
\item \textbf{Two literal forms of \S\ref{sec:p3} are superseded here}, and are corrected here
      rather than there: $r_0=\tfrac12c_1u^{p}$ with "$p\ge1$" (hole~E: $p=0$ and the constant is
      $\beta$-independent), and $\beta_1=\min(1/48,c/N)$ (the scaling is $cN$). Both are
      bookkeeping corrections; neither moves an inequality's direction.
\item \textbf{What is not obtained.} The continuum limit $a\to0$; transport to
      $\beta\approx2.5$; the infinite-volume atom $f_{1,\infty}^2$ and the survival of the atom
      against the continuum threshold at $L=\infty$. \textbf{Clay is not touched: $0\%$.}
\end{itemize}

\noindent\textbf{One reading note.} Where this printing's abstract lists the strong-coupling
result as "$I1$, not $I2$", that records the state of the dossier before the five holes were
closed; nothing in this chapter retracts the abstract's other claims, and nothing in this chapter
enlarges the window or the grade that the abstract states.

\endgroup
